\documentclass[blockverticalspace=3cm]{tikzposter}
%INSTALL

%avoids a warning
%\usepackage[log-declarations=false]{xparse}
\usepackage{xparse}
%Should be inputed before xunicode, someone says. And put everything font-related before fontspec. But I don’t need it anyway.
% \usepackage{amssymb} %for e.g. succcurlyeq.
\usepackage{fontspec} %font selecting commands
%Call amssymb before amsmath or any other amslatex style files (amsthm, ...).
\usepackage{xunicode}
%warn about missing characters
\tracinglostchars=2

%REDAC
\usepackage[hang,flushmargin]{footmisc}%The 'hang' option flushes the footnote marker to the left margin of the page, while the 'flushmargin' option flushes the text as well.
\usepackage{booktabs}
\usepackage{hyphenat}
\usepackage{calc}

\usepackage{mathtools} %load this before babel!
	\mathtoolsset{showonlyrefs,showmanualtags}

\usepackage[super]{nth}
\usepackage{listings} %typeset source code listings
	\lstset{language=XML,tabsize=2,literate={"}{{\tt"}}1,captionpos=b}
\usepackage[nolist,printonlyused]{acronym}%when using smaller, we get: relsize Warning: Failed to get list of defined font sizes.
\usepackage{xspace}
\usepackage[textsize=small]{todonotes}
\usepackage[pdfpagelabels=false]{hyperref}
%pdfusetitle does not work!
%breaklinks makes links on multiple lines into different PDF links to the same target.
%colorlinks (false): Colors the text of links and anchors. The colors chosen depend on the the type of link. In spite of colored boxes, the colored text remains when printing.
%linkcolor=black: this leaves other links in colors, e.g. refs in green, don't print well.
%pdfborder (0 0 1, set to 0 0 0 if colorlinks): width of PDF link border
%hidelinks
\hypersetup{breaklinks,bookmarksopen}
% hyperref doc says: Package bookmark replaces hyperref’s bookmark organization by a new algorithm (...) Therefore I recommend using this package.
\usepackage{bookmark}
\usepackage[capitalise,noabbrev]{cleveref}

% center floats by default, but do not use with float
% \usepackage{floatrow}
% \makeatletter
% \g@addto@macro\@floatboxreset\centering
% \makeatother
\usepackage{enumitem} %follow enumerate by a string saying how to display enumeration
\usepackage{ragged2e} %new com­mands \Cen­ter­ing, \RaggedLeft, and \RaggedRight and new en­vi­ron­ments Cen­ter, FlushLeft, and FlushRight, which set ragged text and are eas­ily con­fig­urable to al­low hy­phen­ation (the cor­re­spond­ing com­mands in LaTeX, all of whose names are lower-case, pre­vent hy­phen­ation al­to­gether). 
\usepackage{siunitx} %[expproduct=tighttimes, decimalsymbol=comma]
\usepackage{braket} %for \Set
\usepackage{natbib}
\usepackage{doi}

\usepackage{amsmath,amsthm}
% \usepackage{amsfonts} %not required?
% \usepackage{dsfont} %for what?
%unicode-math overwrites the following commands from the mathtools package: \dblcolon, \coloneqq, \Coloneqq, \eqqcolon. Using the other colon-like commands from mathtools will lead to inconsistencies. Plus, Using \overbracket and \underbracke from mathtools package. Use \Uoverbracket and \Uunderbracke for original unicode-math definition.
%use exclusively \mathbf and choose math bold style below.
\usepackage[warnings-off={mathtools-colon, mathtools-overbracket}, bold-style=ISO]{unicode-math}

%IJCAI demands Adobe’s Times Roman. The example paper actually uses Nimbus Roman No 9 L (very close); with CM for the sans font. I use texgyretermes, a modern version of NR9L and texgyrecursor for the monospaced version. No instruction provided for math font. I use the default (latin modern).
%\setsansfont[%
%	BoldFont = texgyreheros-bold.otf,%
%	ItalicFont = texgyreheros-italic.otf,%
%	BoldItalicFont = texgyreheros-bolditalic.otf,%
%	Mapping=tex-text% to turn "--" into dashes, useful for bibtex%%
%]{texgyreheros-regular.otf}

\newfontfamily\texgyretermesfamily[
	BoldFont = texgyretermes-bold.otf,%
	ItalicFont = texgyretermes-italic.otf,%
	BoldItalicFont = texgyretermes-bolditalic.otf,%
	Mapping=tex-text% to turn "--" into dashes, useful for bibtex%%
]{texgyretermes-regular.otf}

\newfontfamily\xitsfamily[
	BoldFont = xits-bold.otf,%
	ItalicFont = xits-italic.otf,%
	BoldItalicFont = xits-bolditalic.otf,%
	Mapping=tex-text% to turn "--" into dashes, useful for bibtex%%
]{xits-regular.otf}

%\setmonofont[%
%       Fractions=On,
%       BoldFont = texgyrecursor-bold.otf,%
%       ItalicFont = texgyrecursor-italic.otf,%
%       BoldItalicFont = texgyrecursor-bolditalic.otf,%
%       Mapping=tex-text% to turn "--" into dashes, useful for bibtex%%
%]{texgyrecursor-regular.otf}

%\setmainfont[%
%       Fractions=On,
%       BoldFont = texgyretermes-bold.otf,%
%       ItalicFont = texgyretermes-italic.otf,%
%       BoldItalicFont = texgyretermes-bolditalic.otf,%
%       Mapping=tex-text% to turn "--" into dashes, useful for bibtex%%
%]{texgyretermes-regular.otf}

%NOT loading the font (which should result in latinmodern-math being automatically loaded) solves the spacing problem with some subscripts.
%\setmathfont{latinmodern-math.otf}

% setminus is missing in lmmath font and in texgyrepagella-math
\AtBeginDocument{\renewcommand{\setminus}{\smallsetminus}}
%\setmathfont[range={"29F5}]{xits-math.otf}
%∉
%\setmathfont[range={"2209}]{texgyrepagella-math.otf}
%mapsto
%\setmathfont[range={"21A6}]{texgyrepagella-math.otf}
%✗
%\setmathfont[range={"2717}]{texgyrepagella-regular.otf}
%✓
%\setmathfont[range={"2713}]{texgyrepagella-math.otf}
%\vert, symbol code found in unimath-symbols.pdf (also used for \lvert and \rvert)
%\setmathfont[range={"007C}]{texgyrepagella-math.otf}

%GRAPHICS
\usepackage{pgf}
\usepackage{pgfplots}
\pgfplotsset{compat=1.14}
	\usetikzlibrary{matrix,fit,plotmarks,calc,trees,shapes.geometric,positioning,plothandlers}
\usepackage{graphicx}

\graphicspath{{graphics/},{graphics-dm/}}
\DeclareGraphicsExtensions{.pdf}

%HACKING
\usepackage{printlen}
\uselengthunit{mm}
% 	\newlength{\templ}% or LenTemp?
% 	\setlength{\templ}{6 pt}
% 	\printlength{\templ}
\usepackage{etoolbox} %for addtocmd command
\usepackage{scrhack}% load at end. Corrects a bug in float package, which is outdated but might be used by other packages
\usepackage{xltxtra} %somebody said that this is loaded by fontspec, but does not seem correct: if not loaded explicitly, does not appear in the log and \showhyphens is not corrected.

%Beamer-specific
%\usepackage{todonotes}%UNNECESSARY
%\usepackage{natbib}%REMOVE?
%ADD
% \usepackage{appendixnumberbeamer}
% \newcommand{\citep}{\cite}
% \setbeamersize{text margin left=0.1cm, text margin right=0.1cm} 
% \setbeamertemplate{navigation symbols}{} 
% to solve: this theme produces lots of warnings!
% \usetheme{BrusselsBelgium}
% \usefonttheme{professionalfonts}


\newcommand{\R}{ℝ}
\newcommand{\N}{ℕ}
\newcommand{\Z}{ℤ}
\newcommand{\card}[1]{\lvert{#1}\rvert}
\newcommand{\powerset}[1]{\mathcal{P}(#1)}
\newcommand{\suchthat}{\;\ifnum\currentgrouptype=16 \middle\fi|\;}

\AtBeginDocument{%
	\renewcommand{\epsilon}{\varepsilon}
}

% with amssymb, but I don’t want to use amssymb just for that.
% \newcommand{\restr}[2]{{#1}_{\restriction #2}}
%\newcommand{\restr}[2]{{#1\upharpoonright}_{#2}}
\newcommand{\restr}[2]{{#1|}_{#2}}%sometimes typed out incorrectly within \set.
%\newcommand{\restr}[2]{{#1}_{\vert #2}}%\vert errors when used within \Set and is typed out incorrectly within \set.
\DeclareMathOperator*{\argmax}{arg\,max}


%Voting
\newcommand{\allalts}{\mathcal{A}}
%\newcommand{\allaltsexcept}{\allalts^-}
\newcommand{\allaltsexcept}{\allalts \setminus \set{z}}
\newcommand{\alts}{A}
%the alts on which a specific profile is defined
%\newcommand{\thesealts}[1][\prof]{\alts_{#1}}
\newcommand{\thesealts}{\allalts}
\newcommand{\allvoters}{\mathcal{N}}
\newcommand{\allsystems}{\mathcal{G}}
\newcommand{\prof}{\mathbf{R}}
\newcommand{\tprof}{\prof^*}
\newcommand{\allprofs}{\mathbfcal{R}}
\newcommand{\zerov}{\mathbf{0}}
%all Hamiltonian cycles that need to be considered
\newcommand{\allH}{\mathcal{S}}

\newcommand{\pbasic}[1]{\prof^{#1}_\epsilon}
\newcommand{\pelem}[1]{\prof^{#1}_e}
\newcommand{\pcycl}[1]{\prof^{#1}_c}
\newcommand{\pcycllong}[1]{\prof^{#1}_{cl}}
\newcommand{\pinv}[1][\prof]{\overline{#1}}
\newcommand{\dmap}{{\xitsfamily δ}}
%powerset without zero
\newcommand{\powersetz}[1]{\mathcal{P}_\emptyset(#1)}

%logic atom
\DeclareDocumentCommand{\lato}{ O{\prof} O{\alts} }{[#1 \!\mapsto\! #2]} % UE: added 2 x \! to make it tighter
%logic atom in
\newcommand{\tightoverset}[2]{%
  \mathop{#2}\limits^{\vbox to -.5ex{\kern-0.9ex\hbox{$#1$}\vss}}}
\DeclareDocumentCommand{\latoin}{ O{\prof} O{\alpha} }{[#1 \tightoverset{\in}{⟼} #2]}
\newcommand{\alllang}{\mathcal{L}}
\newcommand{\ltru}{\texttt{T}}
\newcommand{\lfal}{\texttt{F}}
\newcommand{\laxiom}[1]{{\texgyretermesfamily\textsc{#1}}}
\newcommand{\betatr}{\hat{\beta}}
\newcommand{\betasp}{\boxed{β}}
\newcommand{\deltasp}{\boxed{δ}}
\newcommand{\deltacsp}{\boxed{δ\text{-cycl}}}
\newcommand{\kernel}[1]{\mathcal{K}(#1)}
%\DeclareMathOperator{\increas}{inc}
\newcommand{\increas}{\ordaltsrestr{\allaltsexcept}}
\newcommand{\ordaltsrestr}[1]{\restr{{\ordalts}}{#1}}
\newcommand{\ordalts}{\succ}
\newcommand{\ordaltsrev}{\prec}
\newcommand{\ordaltsreveq}{\preceq}
\newcommand{\smlex}{<_{\text{lex}}}

%ARG TH
\newcommand{\AF}{\mathcal{AF}}
\newcommand{\labelling}{L}
\newcommand{\labin}{\textbf{in}\xspace}
\newcommand{\labout}{\textbf{out}}
\newcommand{\labund}{\textbf{undec}\xspace}
\newcommand{\nonemptyor}[2]{\ifthenelse{\equal{#1}{}}{#2}{#1}}
\newcommand{\gextlab}[2][]{
	\labelling{\mathcal{GE}}_{(#2, \nonemptyor{#1}{\ibeatsr{#2}})}
}

\newcommand{\jext}{\hat{j}}

%MISC
\newcommand{\ibeats}{\vartriangleright}
\newcommand{\lequiv}{\Vvdash}
\newcommand{\weightst}{W^{\,t}}
\newcommand{\inferrule}[2]{
	\ensuremath{
		\begin{array}{c}
			#1\\
			\hline
			#2\\
		\end{array}
	}
}
\newcommand{\inferruled}[3]{
	\ensuremath{
		\begin{array}{c|c}
			\multicolumn{2}{c}{#1}\\
			\hline
			#2 & #3\\
		\end{array}
	}
}

%Sorting
\newcommand{\cats}{\mathcal{C}}
\newcommand{\catssubsets}{2^\cats}
\newcommand{\catgg}{\vartriangleright}
\newcommand{\catll}{\vartriangleleft}
\newcommand{\catleq}{\trianglelefteq}
\newcommand{\catgeq}{\trianglerighteq}
\newcommand{\alttoc}[2][x]{(#1 \xrightarrow{} #2)}
\newcommand{\alttocat}[3]{(#2 \xrightarrow{#1} #3)}
\newcommand{\alttoI}{(x \xrightarrow{} \left[\underline{C_x}, \overline{C_x}\right])}
\newcommand{\alttocatdm}[3][t]{\left(#2 \thinspace \raisebox{-3pt}{$\xrightarrow{#1}$}\thinspace #3\right)}
\newcommand{\alttocatatleast}[2]{\left(#1 \thinspace \raisebox{-3pt}{$\xrightarrow[]{≥}$}\thinspace #2\right)}
\newcommand{\alttocatatmost}[2]{\left(#1 \thinspace \raisebox{-3pt}{$\xrightarrow[]{≤}$}\thinspace #2\right)}

\input{preamble/redac}
\input{preamble/draw}
\input{preamble/acronyms}

\title{Arguing about Voting Rules}
\institute{LAMSADE, Université Paris-Dauphine and ILLC, University of Amsterdam}
\author{Olivier Cailloux and Ulle Endriss}
%\usetheme{Basic}

\begin{document}
\maketitle[titletotopverticalspace=10cm]

\begin{columns}
	\column{0.5}
		\block{Main objects}{
\begin{tikzpicture}[remember picture,overlay]
	\path (current page.north west) ++(41cm, -29cm) node[anchor=north west, inner sep=0] (first) {
		\includegraphics[angle=-90, totalheight=10cm]{LOGO.pdf}
	};
	\path (first.south east) ++(2cm, 2.5cm) node[anchor=south west, inner sep=0] (second) {
		\includegraphics[height=6cm]{logo_dauphine_RI_couleurs_edition-1.eps}
	};
	\path (second.east) ++(2cm, -2mm) node[anchor=west, inner sep=0] (third) {
		\includegraphics[height=3.2cm]{logopslnewpng.png}
	};
	\path (third.east) ++(4.3cm, 0) node[anchor=west, inner sep=0] (fourth) {
		\includegraphics[height=4cm]{illc_cropped.pdf}
	};
	\path (fourth.east) ++(1cm, 0) node[anchor=west, inner sep=0] {
	        \includegraphics[height=4cm]{uva.eps}
	};
\end{tikzpicture}
%
\begin{tikzpicture}[remember picture,overlay]
	\path (current page.south west) ++(41.5cm, -25.5cm) node[anchor=south west, text width=23cm] {
		Olivier Cailloux and Ulle Endriss. \emph{Arguing about Voting Rules}. In Proceedings of the 15th International Conference on Autonomous Agents and Multiagent Systems (AAMAS-2016), IFAAMAS.
	};
\end{tikzpicture}%
%
			\begin{itemize}
				\item $\allalts = \{a, b, c\}$ a set of \emph{alternatives}
				\item $N = \{v_1, v_2, v_3\}$ a set of \emph{voters}
				\item $\prof$ a \emph{profile}: maps voters to linear orders on $\allalts$
				\item \emph{Voting rule}: maps profiles to subsets of alternatives: $\lato[\prof][\alts \subseteq \allalts]$
			\end{itemize}
		}
	
	\column{0.5}
		\block{Some example profiles}{
			\begin{equation}
				\prof =
				\begin{array}{ccc}
					v_1&v_2&v_3\\
					a&a&c\\
					b&b&b\\
					c&c&a\\
				\end{array},\quad
				\prof_1 =
				\begin{array}{c}
					v_1\\
					a\\
					b\\
					c\\
				\end{array},\quad
				\prof_2 =
				\begin{array}{cc}
					v_2&v_3\\
					a&c\\
					b&b\\
					c&a\\
				\end{array} 
			\end{equation}
		}
\end{columns}

\begin{columns}
	\column{0.4}
		\block{Our objective}{
			Scenario:
			\begin{itemize}
				\item Assume a commitee wants to choose a voting rule
				\item Many rules
				\item No best rule
			\end{itemize}
			We want to:
				\begin{itemize}
					\item \emph{automatically} explain the outcome of rules
					\item \emph{based} on their properties
					\item by referring to \emph{examples}
				\end{itemize}
		}
	\column{0.6}
		\block{Example dialogue}{
			Let’s argue that in $\prof$, $a$ should win, using “unanimity”, “symmetry” and “reinforcement”
			\begin{center}
				\begin{tabularx}{0.9\colwidth}{>{\bfseries}p{4cm}X}
					System & Consider election $\prof_1$, involving only voter~$v_1$. Do you agree that $a$, enjoying unanimous support, should win this election? \\
					User & Yes, of course. \\
					System & Now consider election $\prof_2$, involving only voters $v_2$ and $v_3$. Do you agree that, for symmetry reasons, the outcome should be a three-way tie? \\
					User & Yes, that sounds reasonable. \\
					System & Observe that when we combine $\prof_1$ and $\prof_2$, we obtain our election of interest, namely $\prof$. Do you agree that in this combined election, as there was a three-way tie in $\prof_2$, $\prof_1$ should be used to decide the winner? \\
					User & Yes, I do. \\
					System & To summarise, you agree that $a$ should win for~$\prof$.
				\end{tabularx}
			\end{center}
		}
\end{columns}

\begin{columns}
	\column{0.5}
		\block{Language}{
			\begin{itemize}
				\item A general language $\alllang$ to talk about properties of voting rules
				\item Logic-based
				\item Can represent axioms and specific outcomes (such as $\lato$)
				\item Permits \emph{arguments}: proof (in the language) from axioms to outcomes %{\tiny (or other properties)}
			\end{itemize}
			\begin{theorem}[Completeness]
				Take some axioms of a voting rule. Encode in our language as $J$. Suppose some outsome $\lato$ holds for every rule $f$ satisfying those axioms. Then there exists a proof of $\lato$ in our logic that is grounded in $J$.
			\end{theorem}
		}
	
	\column{0.5}
		\block{Example for axioms in $\alllang$}{
			\begin{description}
				\item[\laxiom{Reinf}]: $\lato[\prof_1][\alts_1] ∧ \lato[\prof_2][\alts_2] → \lato[\prof_1 \oplus \prof_2][\alts_1 ∩ \alts_2]$ {\tiny with $\alts_1 ∩ \alts_2 ≠ \emptyset$}
				\item[\laxiom{Canc}]: $\prof$ such that $[\forall (a, b): a \succ b \text{ as often as } b \succ a] → \lato[\prof][\allalts]$
				\item[\laxiom{Anon}]: $\lato → \lato[(\prof \circ \sigma)][\alts]$ {\tiny for any permutation $\sigma$}
			\end{description}
		}
\end{columns}

\begin{columns}
	\column{0.3}
		\block{Arguing for Borda}{
			We provide a concrete algorithm for explaining the outcomes of the Borda voting rule.
			Given any profile $\tprof$:
			\begin{itemize}
				\item Let $W \subseteq \allalts$ be the set of Borda winners
				\item The algorithm generates a proof that $\lato[\tprof][W]$
				\item Grounded in axioms: \laxiom{Reinf}, \laxiom{Canc}, \laxiom{Elem}, \laxiom{Cycl}
			\end{itemize}
		}
	\column{0.7}
		\block{Example: Arguing for Borda}{
			Given $\tprof = \begin{array}{ll} a&c\\b&b\\d&a\\c&d\end{array}$, let’s argue that $\lato[\tprof][\{a, b\}]$
			\innerblock{Proof}{%this creates a 4pt overful hbox message, dunno why
				\begin{tikzpicture}
					\newlength{\BESimpleSep}
					\setlength{\BESimpleSep}{0.12\colwidth}
					\path node (Re ab) {$\begin{array}{ll} a&b\\b&a\\c&d\\d&c\end{array}$};
					\path (Re ab.east) node[anchor=west] {$⟼ \{a, b\}$};
					\path (Re ab.east) ++(-2pt, 0pt) node[anchor=south west] {\expl{\laxiom{Elem}}};
					\path (Re ab.east) ++ (\BESimpleSep, 0) node[anchor=west] (Re abc) {$\begin{array}{ll} a&c\\b&b\\c&a\\d&d\end{array}$};
					\path (Re abc.east) node[anchor=west] {$⟼ \{a, b, c\}$};
					\path (Re abc.east) ++(-2pt, 0pt) node[anchor=south west] {\expl{\laxiom{Elem}}};
					\path (Re abc.east) ++ (\BESimpleSep, 0) node[anchor=west] (Rc adcb) {
						$\begin{array}{llll}
							a&d&c&b\\
							d&c&b&a\\
							c&b&a&d\\
							b&a&d&c
						\end{array}$
					};
					\path (Rc adcb.east) node[anchor=west] {$⟼ \allalts$};
					\path (Rc adcb.east) ++(-2pt, 0pt) node[anchor=south west] {\expl{\laxiom{Cycl}}};
					\path (Rc adcb.east) ++ (\BESimpleSep, 0) node[anchor=west] (Rc abdc) {
						$\begin{array}{llll}
							a&b&d&c\\
							b&d&c&a\\
							d&c&a&b\\
							c&a&b&d
						\end{array}$
					};
					\path (Rc abdc.east) node[anchor=west] {$⟼ \allalts$};
					\path (Rc abdc.east) ++(-2pt, 0pt) node[anchor=south west] {\expl{\laxiom{Cycl}}};
					\path ($(Re ab.south east)!0.05!(Rc abdc.south west)$) ++(0, -3em) node[anchor=north] (Rp) {$\prof'$};
					\path[->, draw] (Re ab.south) -- (Rp.north);
					\path[->, draw] (Re abc.south) -- (Rp.north) node[midway, anchor=south] {\tiny 2};
					\path[->, draw] (Rc adcb.south) -- (Rp.north);
					\path[->, draw] (Rc abdc.south) -- (Rp.north);
					\path (Rp.east) node[anchor=west] {$⟼ \{a, b\}$};
					\path (Rp.west) ++(13em, 0) node[anchor=west] {\expl{\laxiom{Reinf}}};
					\path (Rp.south west) ++(0, -1em) node[anchor=west] (Rbig) {$\prof' \oplus \overline{4\tprof} \oplus 4\tprof ⟼ \{a, b\}$};
					\path (Rbig.west) ++(13em, 0) node[anchor=west] {\expl{$\prof' \oplus \overline{4\tprof} \oplus 4\tprof$ must have the same winners as $\prof'$, by \laxiom{Reinf} and \laxiom{Canc}}};
					\path (Rbig.south west) ++(0, -1em) node[anchor=west] (Rbig again) {$4\tprof ⟼ \{a, b\}$};
					\path (Rbig again.west) ++(13em, 0) node[anchor=west] {\expl{$4\tprof$ must have the same winners as $\prof' \oplus \overline{4 \tprof} \oplus 4\tprof$, because $\prof' \oplus \overline{4\tprof}$ cancels}};
					\path (Rbig again.south west) ++(0, -1em) node[anchor=west] (tprof) {$\tprof ⟼ \{a, b\}$};
					\path (tprof.west) ++(13em, 0) node[anchor=west] {\expl{Oh yeah}};
				\end{tikzpicture}
			}
		}
\end{columns}
\end{document}
